\section{A positive Fourier--Jacobi coefficient}
\label{sec:content}

The constant-term estimate would be useless if the entire theta form had
content tending to zero.  We now exhibit a positive Fourier--Jacobi
coefficient whose valuation is bounded above already at the trivial
tame character.
We then impose the additional finite local controls needed for the
Galois argument, while retaining that coefficient bound.  No average
over a group whose order grows with $q_i$ will occur in this section.

\subsection{The norm-one orbit}

Use the Fourier--Jacobi index represented by $x=e_1\in A$, and let
$T=U(Le_2)\subset U(A)$ be its stabilizer.  At nonsplit finite places
choose the unramified hyperbolic test so that its mixed $x$-lattice is
$A(\cO_L)$; keep this Lagrangian lattice and adjust its complement dually.
This is compatible with the scaled binary lattice of
\eqref{eq:theta-spaces}.  Primitive norm-one vectors form one orbit of
the compact unitary group at such a place.

Choose a fixed tensor vector $\varphi\in\pi^\vee$ and a finite-order
character $d_0$ of $[T]$, unramified at nonsplit finite places, for which
\begin{equation}\label{eq:fixed-stabilizer-period}
 I_T(\varphi,d_0):=\int_{[T]}\varphi(s)d_0(s)\,ds\ne0.
\end{equation}
Here is an existence argument with the local conditions included.
A nonzero compact automorphic vector, spherical at nonsplit places, can
be translated at split places to be nonzero somewhere on $T$.
Determinants at nonsplit places are units realized by the maximal compact,
and strong approximation on $SU(A)$ away from a split place removes the
remaining nonsplit components.  Its restriction to the compact torus
quotient is therefore nonzero.  Since its infinity type is trivial,
Fourier expansion on that quotient supplies a finite-order character
$d_0$ with \eqref{eq:fixed-stabilizer-period}.  Sphericity forces $d_0$
to be unramified at the nonsplit places.  A fixed diagonal translation
at $p$ makes $\varphi_p$ invariant under the integral lower unipotent.
This translation multiplies the torus period by a nonzero scalar, so it
does not destroy \eqref{eq:fixed-stabilizer-period}.

For $d_1,d_2\geq0$ satisfying $d_1+d_2=m-1$, take at each real embedding
the polynomial coefficient
\begin{equation}\label{eq:primitive-monomial}
 [e,B_1]^{d_1}[e,B_2]^{d_2}.
\end{equation}
The brackets denote the corresponding $2\times2$ minors.  At $x=e_1$
these minors select the second row of the respective binary columns.

\begin{lemma}\label{lem:one-index-unfolding}
 There are finite tests satisfying the $p$- and $q_i$-requirements of
 Section~\ref{sec:theta} for which the coefficient
 \eqref{eq:primitive-monomial}, at index $(e_1,e_1)$, unfolds into the
 product of a fixed row-one Heisenberg theta function and
 \begin{equation}\label{eq:row-two-unfolding}
  \int_{[T]}\varphi(s)\eta^{-1}(s)
       \vartheta_{1,d_1}(sb_1)\vartheta_{2,d_2}(sb_2)\,ds.
 \end{equation}
 Here $b_1,b_2\in U(1)(\A_F)$ translate the binary lines, and known
 splitting-character factors in $b_j$ have been removed.  At $q_i$ the
 only additional row-two condition is a unit cut in the first binary
 column.  The unfolding introduces a fixed nonzero volume factor and
 no factor $q_i-1$.
\end{lemma}
\begin{proof}
 At a fixed split place away from $p$, use small support transverse to
 the orbit of $e_1$ in the mixed polarization, leaving arbitrary desired
 unary tests on the remaining coordinates.  Unfolding the rational
 norm-one orbit then cuts the finite integration to its stabilizer
 $T$.  At a nonsplit place the single compact orbit of primitive
 norm-one vectors gives the same cut, with all remaining functions
 spherical.

 At $q=q_i$, integral vector--covector pairs pairing to $1$ are one
 $\GL_2(\cO_q)$-orbit.  The pivot test
 \eqref{eq:q-pivot-test} therefore effects the same cut.  At augmentation
 its determinant character is trivial.  Setting the position vector
 equal to $e_1$ makes invertibility of the head precisely the condition
 that the second entry of $Y_{B_1}$ is a unit; the other remaining
 entries are integral.  The orbit volume is the unmodified ratio
 $\vol(K_{A,q})/\vol(T(\cO_q))$, so it does not introduce the residue
 unit-group order.

 At $p$, take $Y_e$ in a small fixed neighborhood of $e_1$,
 $Y_{B_1}$ integral with any prescribed locally constant test on its
 second entry supported on units, and $Y_{B_2},Y_f$ integral.  The dual
 coordinate of $x$ is integral and pairs to $1$ with $Y_e$.
 Such pairs form one orbit under a sufficiently deep source congruence
 subgroup enlarged by the integral lower unipotent.  The latter fixes
 $\varphi_p$ by its choice above, and its transpose on the covector
 test leaves the relevant second entry unchanged.  Thus the row-two
 $B_1$ test can have arbitrary fixed locally constant unit level while
 the row-two $B_2$ test is the integer test.  The depth of $Y_e$ and all
 other $p$-levels can be fixed once and for all.  These tests still have
 the form \eqref{eq:p-pivot-test}.

 At infinity, rotating $x$ introduces exactly the determinant
 multiplier already present in the source type.  The minor coefficient
 therefore becomes the specified second-row powers.  The Heisenberg
 translation acts only on the first row; the second row gives
 \eqref{eq:row-two-unfolding}.  On a line, the two torus factors act by
 their product, so translating the binary argument by $b_j$ is the
 same as $s\mapsto sb_j$, up to its explicitly removable splitting
 character.  This proves the factorization.
\end{proof}

The row-one function in this lemma is nonzero at some prime-to-$p$
Jacobi torsion for every one of the finitely many torus representatives
that will be needed.  Indeed, the Heisenberg realization with nontrivial
central character is injective by its Fourier expansion (or by the
irreducibility of that representation).  More concretely, with finite
components fixed, finite-adelic approximation finds a rational position
where the translated finite Schwartz function is nonzero.  Its real
Gaussian is nonzero, and modulation in the real Heisenberg variable
separates the distinct rational positions.  Prime-to-$p$ torsion
arguments, with compatible finite Heisenberg shifts if necessary, are
dense; they therefore detect this nonzero function.  Move each chosen
torsion to zero by a rational unipotent and use the translated split
differential frames.  At $p$ the parameters can be integral and preserve
the fixed deep level.  These moves affect only row one.  Their normalized
values form a fixed finite list of nonzero constants, costing only a
fixed denominator and valuation.

\subsection{The two unary characters and their separate root numbers}

Split the action of $T$ in \eqref{eq:row-two-unfolding} into the two
line oscillators.  Choose inverse splitting characters for the binary
columns, each restricting to $\omega_{L/F}$ on $\A_F^\times$, with
opposite infinity exponents $\pm\frac12$ in unitary base-change
normalization, that is, $z\mapsto(z/|z|)^{\pm1}$.  They may be taken
unramified at $p$ and at every nonsplit finite place, with all additional
conductor supported at split places.  To construct them, prescribe the
restriction, infinity types, and unit conditions in the ray class
construction.  A possible principal relation has ratio to its
conjugate a unit root of unity.  A sufficiently deep split conductor
forces that ratio to be $1$, reducing the compatibility condition to
the prescribed one on $F^\times$.  Extend over the remaining finite ray
class quotient and take the inverse for the other column.

Use compatible tensor splittings.  Their product is trivial, as is
required by the binary action of Section~\ref{sec:theta}.  Alternatively,
a discrepancy between such splittings is a fixed finite-order character:
its infinity type is trivial by the vacuum formula, and its nonsplit
finite components are spherical by the matched lattices.  It can be
absorbed into the finite character product below, in its first slot at
$p$.  The hyperbolic factor uses the trivial splitting and contributes
no stabilizer character.

The two $T_\infty$-weights are consequently parallel $d_j+c_j$, where
$c_j$ are fixed integers and $c_1+c_2=1$.  Choose characters $\xi_j$ of
these weights with
\begin{equation}\label{eq:unary-character-product}
 \xi_1\xi_2=\cC^m d_0,
\end{equation}
unramified at nonsplit places, and with $\xi_2$ unramified at $p$.
They are powers of $\cC$ times fixed characters; the first slot carries
any required fixed $p$-character.  The same split-conductor construction
makes these choices possible.

Averaging the translates in \eqref{eq:row-two-unfolding} against
$\xi_j(b_j)^{-1}$ separates the variables.  At the trivial tame
character the result is
\begin{equation}\label{eq:unary-period-factorization}
 I_T(\varphi,d_0)\,T_1(d_1)\,T_2(d_2),\qquad
 T_j(d_j)=\int_{[U(1)]}\vartheta_{j,d_j}(s)\xi_j(s)^{-1}\,ds,
\end{equation}
with the nonzero row-one factor understood.  For the moment omit the
extra unit cut at $q_i$.

The two degrees must eventually satisfy a joint condition:
$m_i=d_{1i}+d_{2i}+1$ tends to zero $p$-adically.  Their $p$-adic
limits will therefore lie on the line $d_1+d_2=-1$, although the
classical approximating degrees are positive and tend to infinity.
We will obtain a nonzero analytic period for each line oscillator and
then choose a point on this compatible pair of degree discs avoiding
their zeros.

The discs can be fixed before choosing the auxiliary finite-order
characters in the next lemma.  Represent the torus classes integrally
at $p$, translate only their finite tests, and evaluate at fixed CM
frames.  The factors varying with degree are powers of the fixed
$\cC$, the fixed unit period and coordinate factors, and the unit
moments of Lemma~\ref{lem:unary-moments}.  The values of $\cC$ on
these representatives lie in the units of a fixed finite extension.
A new prime-to-$p$ finite test changes its fixed coefficients, but not
the width required for those variable powers.  Thus, for every fixed
choice of the finite characters and tests, sufficiently narrow fixed
residue discs give convergent analytic functions $\mathcal T_j(s)$
for the unit-cut periods.  The lemma supplies a nonzero classical
value for each function and makes the two discs compatible with
$d_1+d_2=-1$.

\begin{lemma}\label{lem:simultaneous-unary-nonvanishing}
 The fixed finite-order parts of $\xi_1,\xi_2$ can be chosen, preserving
 \eqref{eq:unary-character-product}, so that suitable permitted local
 tests give $T_j(d_{j0})\ne0$ at positive classical cut weights
 $d_{10},d_{20}$.  The weights can be chosen so that
 $d_{10}+d_{20}+1$ is divisible by every prescribed torsion order and
 a prescribed sufficiently large fixed power of $p$.
\end{lemma}
\begin{proof}
 Begin with large weights satisfying the indicated congruences.  The
 local theta Hom for each individual line is nonzero: at a nonsplit
 finite place this is the unramified correspondence for the matched
 self-dual lattice, at infinity it is the chosen Fock degree, and at a
 split place there is no local obstruction.

 We check the global root number separately for each line.  If $\alpha$
 is its $U(1)$-character and $\chi_W$ the corresponding splitting
 character, the standard unitary Hecke character entering its central
 $L$-value is
 \[
  \chi_W^{-1}\alpha_L,\qquad
  \alpha_L(z)=\alpha(z/z^c).
 \]
 Since $\alpha_L$ is trivial on $\A_F^\times$ and
 $\chi_W|_{\A_F^\times}=\omega_{L/F}$, this character restricts to
 $\omega_{L/F}$.  The dimension-one epsilon dichotomy identifies its
 local epsilon factor, for the chosen trace-zero element and additive
 character, with the product of the two local line invariants
 \cite[Theorem~3.5]{BoradeEtAl2025}.  The product formula for these
 invariants is $1$ for each coherent global pair of lines.  Thus each
 of the two global root numbers is $+1$; it is not only their product
 that is known to be positive.

 At infinity the pullback of $\xi_j$ has exponent $2(d_j+c_j)$,
 modified by $\pm1$ by the splitting.  After the arithmetic norm-half
 shift, and if necessary conjugating the CM orientation using
 self-duality, its critical algebraic character has type
 \begin{equation}\label{eq:unary-critical-type}
  (a_j+1)\Sigma-a_j\Sigma^c,
  \qquad a_j=d_j+O(1)\geq0,
 \end{equation}
 and restriction $\omega_{L/F}|\cdot|_{\A_F}$.  Its conductor has only
 split prime factors.

 For the simultaneous nonvanishing, use the characteristic-zero
 theorem of Burungale--He--Tian--Ye
 \cite[Theorem~1.1 and Corollary~4.20]{BHTY2026}.
 Write $u_j=\chi_{W,j}^{-1}(\xi_j)_L$ for the unitary standard
 character of the $j$th line and put
 \[
  \lambda_j=u_j|\cdot|_{\A_L}^{-1/2}.
 \]
 This is the opposite norm-half shift from the arithmetic character
 in \eqref{eq:unary-critical-type}.  Since the idelic norm on $L$
 restricts to the square of that on $F$,
 \[
  \lambda_j|_{\A_F^\times}
    =\omega_{L/F}|\cdot|_{\A_F}^{-1},\qquad
  L(1,\lambda_j\nu)=L(1/2,u_j\nu).
 \]
 Thus $\lambda_j$ has the self-duality convention of that theorem.
 For the large weights already chosen, its infinity type has the
 required form after choosing the corresponding CM orientation; the
 theorem allows an arbitrary CM type, including an induced one.

 Choose a degree-one prime $\mathfrak l$ of $F$ above an odd rational
 prime $\ell$, split in $L$ and outside all fixed conductors and $p$.
 Let $\Gamma_{\mathfrak l}$ be the Galois group of the maximal
 $\Z_\ell$-free anticyclotomic extension of $L$ unramified outside
 $\mathfrak l$; it is isomorphic to $\Z_\ell$.
 The two base signs are separately $+1$ by the preceding local
 calculation.  At this split prime they remain $+1$ for all sufficiently
 ramified finite-order characters of $\Gamma_{\mathfrak l}$
 \cite[Lemma~4.12 and Section~4.4.3]{BHTY2026}.
 The nonvanishing theorem therefore gives finite exceptional sets for
 each of $L(1/2,u_1\nu)$ and $L(1/2,u_2\nu^{-1})$.
 Choose $\nu$ outside their union.

 Because $\ell$ is odd, $\nu$ has a unique square root among the
 $\ell$-power-order characters.  Restrict $\nu^{1/2}$ to
 $\U(1)(\A_F)$ and call the resulting character $\tau$.
 Anticyclotomicity gives
 \[
  \tau(z/z^c)=\nu^{1/2}(z/z^c)=\nu(z).
 \]
 Replacing $\xi_1,\xi_2$ by $\xi_1\tau,\xi_2\tau^{-1}$
 preserves Equation~\ref{eq:unary-character-product} and gives the
 two nonzero central values just selected.  This adds conductor only at
 the split prime $\mathfrak l$, so all prescribed local conditions are
 preserved.  Only characteristic-zero nonvanishing is used here.

 The equal-rank line theta norm formula, equivalently the global
 line-lift criterion
 \cite[Theorem~4.2 and Proposition~4.3]{BoradeEtAl2025}, now supplies nonzero unary periods
 with suitable local tests.  At $p$ use the appropriate unit-character
 test for row-two $B_1$ and the unit cut of the unramified test for
 row-two $B_2$; their split local Mellin integrals are nonzero.
 Keep the spherical tests at nonsplit places.  Compact line groups
 and the nontrivial infinity powers give the cuspidal situation with
 no occurrence from dimension zero.  The nonzero output on the other
 line is a character, so it is already nonzero at identity: translating
 that argument factors through the input torus, up to the known
 splitting scalar.  This proves the assertion at the chosen cut
 weights.  Fix this auxiliary finite-order choice permanently.  Its
 conductor and every additional local denominator are now fixed;
 enlarge the finite set in Proposition~\ref{prop:tame-primes} before
 the final moving-prime sequence is chosen.
\end{proof}

\subsection{Interpolation and the moving-prime unit cut}

\begin{lemma}\label{lem:bounded-uncontrolled-content}
 Before imposing the additional fixed-place controls, there are
 positive $d_{1i},d_{2i}$, with $m_i=d_{1i}+d_{2i}+1$ tending to zero
 $p$-adically and to infinity archimedeanly, for which the theta forms
 have a nonconstant Fourier--Jacobi coefficient of valuation bounded
 above at augmentation.  The choices can simultaneously satisfy the
 constant-term precision of Proposition~\ref{prop:constant-term-small}.
\end{lemma}
\begin{proof}
 Apply Lemma~\ref{lem:unary-moments} to the two unary periods.  For the
 actual $B_2$ integer test, positive degrees tending to infinity pick
 out its unit cut.  Thus the analytic limiting function specializes at
 $d_{20}$ to the nonzero cut value obtained in
 Lemma~\ref{lem:simultaneous-unary-nonvanishing}; the same is true for
 $B_1$.  Hence neither $\mathcal T_j$ is identically zero on its
 chosen disc.  This is the reason for testing nonvanishing with unit cuts,
 while retaining the integer test in the actual third column of
 \eqref{eq:p-pivot-test}.

 By the congruence imposed on $d_{10}+d_{20}+1$, the discs can be
 arranged so that $s$ in the first disc implies $-1-s$ in the second.
 The analytic functions $\mathcal T_1(s)$ and
 $\mathcal T_2(-1-s)$ are nonzero functions.  Choose a point $a$ avoiding
 their zero sets.  We may take positive integer sequences
 $d_{1i}\to a$ and $d_{2i}\to-1-a$, both tending to infinity, on the
 required torsion branches.  Their sum plus one then tends to zero.
 The two period values stay bounded away from zero.

 Restore the unit cut above $q_i$.  Subtracting the uniformizer
 translate of the relevant unary Schwartz function from its integer
 test and using torus equivariance expresses its effect by finitely
 many factors
 \begin{equation}\label{eq:content-moving-euler}
  1-\alpha_{iv}\cC(\Frob_{iv})^{\pm d_1}.
 \end{equation}
 Here $\alpha_{iv}$ consists of residual-cardinality powers, fixed
 finite characters, splitting characters, and the fixed weight shifts.
 Its limit is a nonzero $p$-unit; the dependence through
 $d_2=m-1-d_1$ is incorporated in the sign and these shifts.
 Pass to sublimits for the bounded-degree roots occurring in these
 normalizations.  Proposition~\ref{prop:tame-primes} gives
 $\cC(\Frob_{iv})\to\cC(\gamma)$, which is a $p$-unit because
 $\cC$ is a continuous $p$-adic Galois character, and is not a root of unity.
 On a fixed weight branch its unit logarithm is consequently nonzero.
 None of the limiting analytic factors
 \eqref{eq:content-moving-euler} is identically zero.  Refine the
 choice of $a$ to avoid their finitely many zero sets as well.
 The product of the two unary periods, all these local factors,
 the fixed period \eqref{eq:fixed-stabilizer-period}, and the chosen
 row-one evaluations is then bounded away from zero.  The remaining
 $m$-dependent normalization factors in
 \eqref{eq:primitive-monomial} are only the unit powers already
 treated in the mixed expansion.

 There is no growing averaging denominator in turning this product
 back into a coefficient.  The auxiliary finite-order choices and
 the torus levels have been fixed; the $q_i$ unit tests are invariant
 under local torus units.  Thus the averages over the $b_j$ in
 \eqref{eq:unary-period-factorization} use a fixed finite set of
 representatives, chosen away from the varying $q_i$.  If every one
 of the corresponding coefficient evaluations had valuation tending
 to infinity, so would their fixed finite weighted sum, contrary to
 the product just obtained.  Hence one of finitely many evaluations
 of the positive-index coefficient has valuation bounded above.
 Passing to the filter chooses one option if needed.

 Finally choose the positive integers $d_{1i},d_{2i}$ sufficiently
 close to their limits at each stage, while keeping them arbitrarily
 large.  Their sum plus one, $m_i$, can be as close to zero as required by
 the fixed-$i$ constant-term estimate.  Simultaneous congruences and
 lower bounds on the positive integers cause no conflict.  This
 proves the assertion with the required constant-term precision.
\end{proof}

\subsection{Fixed-place replacements with bounded denominators}

We now prescribe the remaining levels at split places outside $p$.
The comparison point is $\eta=1$, where
\eqref{eq:theta-eigenparameters} consists of the elliptic block and the
final pair $|\cdot|^{-1/2},|\cdot|^{1/2}$.  The separation used here
excludes a ratio $Q_v^n$, $n\in\Z$, between a Frobenius root of the
elliptic block and one of the final pair; $Q_v$ is the residue
cardinality.  This is the separation applied in
Lemma~\ref{lem:binary-inertia}.
At a place of potentially good reduction, or whenever the two blocks of
\eqref{eq:theta-eigenparameters} have no norm-power linkage at the
comparison point, use an upper $(2,2)$ level: deep first Levi and
opposite radical, full last Levi and upper radical.  This includes a
ramified quadratic Steinberg twist when the Frobenius lift is chosen
with its quadratic minus sign.  Use open-pivot tests with
$Y_{[1,2]}\in\GL_2(\cO)$, nontrivially paired with the local source
vector, and the remaining columns integral.  At an unramified Steinberg
twist where linkage may occur, use a maximal $(3,1)$ parahoric and make
no additional Hecke refinement there.

\begin{lemma}\label{lem:fixed-place-spanning}
 At each such fixed split place, the tests just prescribed and their
 translates span the tests used in
 Lemma~\ref{lem:bounded-uncontrolled-content} in theta coinvariants,
 after inverting a Laurent polynomial nonzero at weight zero.  For
 weights tending to zero, this spanning uses finitely many translates
 and a common bounded denominator.  All replacements can be made
 independently while leaving the $p$- and $q_i$-tests unchanged.
\end{lemma}
\begin{proof}
 Work in the split local theta coinvariants over the Laurent ring of
 the unramified determinant deformation, localized near its central
 value.  Filter locally constant Schwartz functions on $M_{2,4}$ by
 rank.  The rank-zero contribution would require a character quotient
 of the generic source representation, which is impossible.
 For a rank-one image fix its first covector.  The stabilizer Levi
 $\diag(1,d)$ would require exponent $|d|^{-1}$ on the raw lower
 Jacquet module of $\pi^\vee$: the Weil factor is $|d|^2$ and the
 vector-orbit Jacobian is $|d|$.  In normalized Jacquet conventions
 the required exponent is $|d|^{-3/2}$, which is excluded by
 temperedness of the elliptic source representation.  For a
 supercuspidal source the relevant Jacquet module is zero; for a
 tempered principal or special source its exponents have the usual
 tempered bounds and do not equal $-3/2$.

 These are calculations over the deformation ring, not merely at
 isolated characters.  Use restriction of locally constant functions
 for the rank filtration and integration along the vector orbit, or
 smooth compact-induction coinvariants, on compact open charts of
 the row/image parameters.  The difference between the required Levi
 character and each of the finitely many actual Jacquet characters
 is invertible after localization at the comparison point.  Thus
 the lower-rank coinvariants vanish in that localization, and the
 open rank-two part surjects onto all coinvariants.

 On the open rank-two orbit, direct orbit induction identifies the
 coinvariants with normalized parabolic induction of the
 determinant-twisted source and the trivial binary representation.
 At weight zero this induction is irreducible by irreducibility of
 unitary parabolic products for general linear groups
 \cite{Tadic1986}.  The pivot test is nonzero there, since its head
 can be paired with a nonzero local matrix coefficient.  Its
 translates consequently generate.  In the Steinberg case there are
 nonzero $(3,1)$-parahoric invariants: in the induced model place the
 two Steinberg positions on different sides of the single-step flag.
 Their finite-Weyl sign condition then imposes no forbidden merging
 of those positions.  Project an open test nontrivially to these
 invariants; its translates also generate the central
 specialization.

 Only finitely many original test vectors at one fixed level must be
 spanned.  In that finite-dimensional induced model, select a finite
 collection of the permitted translates whose spanning matrix has a
 minor $f$ with $f(1)\ne0$.  Cramer's rule expresses the needed
 vectors after inverting $f$.  The coefficients are rational
 functions of the unramified parameter.  Since
 $\eta_m\to1$ and $\nu_i$ is locally trivial at this fixed place,
 $f(\eta_m)\to f(1)\ne0$.  All these coefficients therefore have
 bounded denominator.  There are only finitely many fixed places,
 so tensoring the finite expressions gives one bound.  No inverse
 involving the varying place $q_i$ has entered this argument.
\end{proof}

If all the resulting forms with the prescribed controls had content
tending to zero, then their finitely many fixed prime-to-$p$ translates
would too.  Such translations preserve ordinary integrality: they are
prime-to-$p$ isogeny correspondences with unit differential comparisons,
and the finitely many fixed averaging factors cost at most a fixed
constant.  The bounded expression in
Lemma~\ref{lem:fixed-place-spanning} would force the original form of
Lemma~\ref{lem:bounded-uncontrolled-content} to have content tending to
zero, a contradiction.  At least one controlled choice therefore retains
the bounded-content coefficient.  A filter subsequence fixes the choice
if necessary.  If a finer neat level is needed, pass to a fixed-index
neat cover at one fixed auxiliary place and descend by its fixed
averaging factor; this changes only the uniform constant.

\subsection{Positive Hecke controls on the final binary block}

At each upper $(2,2)$ level, including the moving $q_i$ level, retain the
strictly expanding center operator $U_2$ and the degree-one spherical
operator of the last $\GL_2$ in the normalized Jacquet module.  They act
on our test with the final allocation in
\eqref{eq:theta-eigenparameters}, namely the two characters
$|\cdot|^{-1/2},|\cdot|^{1/2}$ alone.  In particular the center
expansion value is invertible.

For completeness this assertion follows at the level of the tests.
Represent double cosets of the expanding monoid in the last block and
lift them with the upper shifts.  The invertible head determines exactly
one surviving upper shift for each last-block coset.  For both
$\varpi^{-1}1_2$ and $\diag(\varpi^{-1},1)$, the determinant factor in
\eqref{eq:split-weil-action} is exactly the factor that undoes the
Jacquet half-modulus.  After normalized Jacquet passage, the two
integral trailing columns therefore carry the spherical action of the
trivial binary representation.  Its Satake pair is
$(Q^{-1/2},Q^{1/2})$, where $Q$ is the residue cardinality.  The first
block may have the principal congruence level needed at $q_i$; it does
not change this count.

Iwahori decomposition for the block parabolic makes these positive
operators commute.  Their normalizations use only powers of the local
residue characteristic.  Inverting the strictly contracting operator
identifies them with the stated Jacquet operators by Jacquet lifting
\cite{Casselman1995}.  At the fixed places the replacements in
Lemma~\ref{lem:fixed-place-spanning} can be made with these tests from
the outset.  The $p$- and $q_i$-requirements remain unchanged.

\begin{proposition}[The combined theta output]\label{prop:bounded-content}
 Fix the truncation $b_*$ selected in Section~\ref{sec:tame}.  There
 exist positive integers $m_i\to\infty$, tending $p$-adically to zero
 on the prescribed torsion branch, and theta forms $\Theta_i$ of
 weight $(m_i,m_i,1;-1)$ with the following simultaneous properties.
 \begin{enumerate}[label=\textup{(\roman*)}]
 \item Their ordinary monomial-dual lattices have a uniform denominator
       bound, coefficientwise in the full tame group algebra and after
       every fixed truncated substitution used in the argument.  One
       fixed scaling clears this bound.
 \item All constant Fourier--Jacobi coefficients tend to zero through
       degree $b_*-1$ of the tame variable, in compatible ordinary
       differential bases, uniformly at the boundary representatives.
       This is a full-ring estimate in $R_{b_*}$.
 \item At augmentation, a nonconstant Fourier--Jacobi coefficient has
       an integral-frame evaluation of valuation bounded above by a
       constant independent of $i$.  Thus the forms do not acquire
       unbounded common content.
 \item At both distinguished places above $p$, the fixed canonical
       $(2,1,1)$ level is $J$ and the operators $p^{-2}U_2,p^{-2}U_3$
       have value $1$.  At every controlled upper $(2,2)$ place,
       including $q_i$, the expanding center is invertible and the
       last binary spherical data are those of the final pair in
       \eqref{eq:theta-eigenparameters}.  At the specified unramified
       Steinberg places the level is $(3,1)$ parahoric, without an
       additional refinement.  All nonsplit finite places are
       hyperspecial.
 \item At split good spherical places the Hecke parameters are
       \eqref{eq:theta-eigenparameters} with
       $\eta=\cC^{m_i}\nu_i$.
 \end{enumerate}
 All constants may depend on the fixed original data, the fixed local
 tests and $b_*$, but not on $i,q_i,n_i$ or the growing weights.
\end{proposition}
\begin{proof}
 Lemma~\ref{lem:bounded-uncontrolled-content} chooses the degrees and
 obtains the content bound simultaneously with
 Proposition~\ref{prop:constant-term-small}.  The bounded fixed-place
 spanning and the preceding content argument impose all additional
 controls without losing that bound.  Every selected test still
 satisfies the hypotheses of Proposition~\ref{prop:theta-family} and
 the constant-term proof, which allowed arbitrary bounded fixed split
 tests.  The explicit positive-operator calculation gives the last
 binary allocations.  If finitely many choices remain, make the
 constant-term approximation sufficiently accurate for all of them
 and then choose the one with bounded content.  This gives all the
 properties at once.
\end{proof}
